% !PRINTED 491
\weberpage{491}

If we set $u = v$, then

\[
A = 2, \quad B = u^2 + \frac{2}{u^2},
\]
and if we set $u = f(\sqrt{-m})$, $m = 7, 19, 23$, then from \S\,128, (6), (8), \S\,131, (1) we obtain by elimination of $u$ the equations for $B$:

\[
\begin{aligned}
m &= 7, & B - 3 &= 0, \\
m &= 19, & B^2 - 6B^2 + 10B - 6 &= 0, \\ % [sic] printed first term is $B^2$; the sentence below and (7) require $B^3$
m &= 23, & B^3 - 5B^2 + 4B - 1 &= 0.
\end{aligned}
\]

The two last cubic equations are irreducible, since they have no rational root, and from this the required modular equation (3) follows without calculation:

\weq{7}{A - 2 = (B-3)^2\,(B^3 - 6B^2 + 10B - 6)^2\,(B^3 - 5B^2 + 4B - 1).}

In the memoir: \glqq Ein Beitrag zur Transformationstheorie der elliptischen Funktionen mit einer Anwendung auf Zahlentheorie\grqq{}, Math. Annalen 43, 185, I have derived by the same route the transformation equation for the 47th degree of transformation as well.

\begin{center}
\S\,133. \gesperrt{The resolvents of degrees 7 and 11 for the 7th and 11th degree of transformation.}
\end{center}

We have seen in \S\,82 that for the 7th and 11th degree of transformation there exist resolvents of degrees 7 and 11. In the formation of these equations too the theory of complex multiplication can render useful service.

When $n = 7$ or $= 11$, we consider the transformation equation

\[
F_n(u, v) = 0,
\]
whose roots, if we set $u = f(\omega)$, are

\weq{1}{v_\infty = f(n\,\omega), \quad v_c = f\left(\frac{\omega+c}{n}\right)}

where $c$ may be taken as index of $v$ with respect to the modulus $n$, while under the symbol $f$ it must be assumed divisible by 48.

These equations are, according to \S\,73:

\weq{2}{n = 7, \quad v^8 - u^7v^7 + 7u^4v^4 - 8\,uv + u^8 = 0,}

\weq{3}{n = 11, \quad v^{12} - u^{11}v^{11} + 11\,u^9v^9 - 44\,u^7v^7 + 88\,u^5v^5 - 88\,u^3v^3 + 32\,uv + u^{12} = 0.}
